В ходе выполнения данной курсовой работы были изучены и реализованы два
классических алгоритма сортировки"--- быстрая сортировка и сортировка
слиянием.

В результате проведённого теоретического анализа установлено, что оба
алгоритма имеют одинаковую среднюю временну\'ю сложность $O(n \log n)$,
однако существенно различаются в поведении при неблагоприятных входных
данных: быстрая сортировка деградирует до $O(n^2)$, тогда как сортировка
слиянием гарантирует $O(n \log n)$ при любом наборе данных.

Эмпирические эксперименты подтвердили теоретические выводы: на случайных
данных быстрая сортировка незначительно быстрее за счёт лучшей
кэш"=локальности, однако уступает по предсказуемости поведения.

На основе проведённого анализа можно сформулировать следующие
практические рекомендации:
\begin{itemize}
    \item если память ограничена, а данные случайны"--- предпочтительна
          быстрая сортировка;
    \item если требуется гарантированное время $O(n \log n)$ или
          стабильность"--- предпочтительна сортировка слиянием;
    \item в стандартных библиотеках (например, \texttt{std::stable\_sort}
          в C++) применяется именно сортировка слиянием~\cite{cppreference}.
\end{itemize}
